\documentclass[12pt,a4paper]{article}

% ------------------------------------------------
% Pacotes
% ------------------------------------------------
\usepackage[utf8]{inputenc}
\usepackage[T1]{fontenc}
\usepackage[brazil]{babel}
\usepackage{graphicx}
\usepackage[table]{xcolor}
\usepackage{array}
\usepackage{tabularx}
\usepackage{geometry}
\usepackage{fancyhdr}
\usepackage{mathptmx} % Times New Roman
\usepackage{enumitem}\usepackage{amsfonts}
\usepackage{amsmath} % Recomendado para o comando \text{}
\usepackage{amssymb}
\usepackage{tikz}
\usetikzlibrary{positioning, shapes.geometric, arrows.meta}

% TikZ Styles for ER Diagrams
\tikzset{
    block/.style={rectangle, draw, thick, minimum width=2cm, minimum height=1cm, fill=white},
    diamond_node/.style={diamond, draw, thick, minimum width=2cm, minimum height=1cm, fill=white, aspect=2},
    attribute/.style={ellipse, draw, thin, minimum width=1.5cm, minimum height=0.6cm, fill=white, font=\scriptsize},
    line/.style={draw, thick},
}
% ------------------------------------------------ 
% Configuração de página 
% ------------------------------------------------
\geometry{
    a4paper,
    left=2cm,
    right=2cm,
    top=2cm,
    bottom=2cm,
    headheight=15pt 
}

% Caixa para resposta discursiva
\newcommand{\resposta}[2]{
\noindent
\framebox[\textwidth]{
\begin{minipage}{0.95\textwidth}
\vspace{0.2cm}
\textbf{RESPOSTA:} \\
#2
\vspace{#1}
\end{minipage}
}
}

\pagestyle{fancy}
\fancyhf{}
\renewcommand{\headrulewidth}{0pt}

\setlength{\parindent}{0pt}

% ------------------------------------------------
% Comandos auxiliares
% ------------------------------------------------
\newcolumntype{L}[1]{>{\raggedright\arraybackslash}p{#1}}
\newcolumntype{C}[1]{>{\centering\arraybackslash}p{#1}}

% ------------------------------------------------
% Documento
% ------------------------------------------------
\begin{document}

% ------------------------------------------------
% Cabeçalho
% ------------------------------------------------
\begin{tabularx}{\textwidth}{@{} L{4cm} X @{}}
    \raisebox{-0.5cm}{\includegraphics[width=4cm]{Universidade_de_Rio_Verde.jpg}} &
    \hfill
    \begin{tabular}{|>{\centering\arraybackslash}m{9.5cm}|>{\raggedright\arraybackslash}p{1.8cm}|}
        \hline
        \cellcolor{black}\color{white}\rule{0pt}{1.2cm}\textbf{GABARITO - AVALIAÇÃO 2} 
        & 
        \scriptsize Nota \tabularnewline
        \hline
    \end{tabular} 
\end{tabularx}
\vspace{0.3cm}

% ------------------------------------------------  
% Tabela principal
% ------------------------------------------------
\begin{tabularx}{\textwidth}{|X|L{2cm}|}
\hline
\textbf{Curso} & \textbf{Campus} \\
\textit{Engenharia de Software} & \textit{Rio Verde} \\
\hline
\multicolumn{2}{|l|}{\textbf{Disciplina}} \\  
\multicolumn{2}{|l|}{ \textit{Programação Orientada a Objetos}} \\
\hline
\textbf{Nome do(a) acadêmico(a)} & \textbf{Valor} \\
\textbf{GABARITO OFICIAL} & 7,0 \\
\hline
\textbf{Professor} & \textbf{Turma} \\
\textit{ Sergio Souza Novak} & \textit{B} \\
\hline
\textbf{Data da Avaliação} & \textbf{Código} \\
\textit{13/05/2026} & \textit{ESW448} \\
\hline
\end{tabularx}
 
\vspace{0.5cm}

\begin{center}
    \rule{\textwidth}{0.4pt}
    \textbf{QUESTÕES TEÓRICAS}
\end{center}

\vspace{0.2cm}

\begin{enumerate}
    \item \textbf{Resposta: B}
    \begin{itemize}
        \item Afirmativa 1: Falsa. Classes abstratas podem conter atributos normalmente.
        \item Afirmativa 2: Verdadeira. Podem conter tanto métodos abstratos quanto concretos.
        \item Afirmativa 3: Falsa. Construtores não podem ser abstratos.
        \item Afirmativa 4: Falsa. Classes abstratas não podem ser instanciadas diretamente.
    \end{itemize}

    \item \textbf{Resposta:}\\
    \resposta{0.5cm}{
    O modificador \texttt{private} restringe o acesso ao atributo \texttt{nome} apenas à própria classe \texttt{Animal}, impedindo o acesso direto pelas subclasses. O modificador \texttt{protected} permite que o atributo \texttt{idade} seja acessado por \texttt{Animal}, suas subclasses (\texttt{Cachorro} e \texttt{Gato}) e outras classes do mesmo pacote. O modificador \texttt{public} torna o método \texttt{emitirSom()} acessível de qualquer lugar. \\
    No método \texttt{exibirDados()}, o atributo \texttt{nome} \textbf{não} pode ser acessado diretamente devido ao encapsulamento (\texttt{private}), enquanto o atributo \texttt{idade} \textbf{pode} ser acessado diretamente por conta da herança e do modificador \texttt{protected}.
    }

    \item \textbf{Resposta:}
    \begin{enumerate}[label=( \quad )]
        \item[\textbf{( V )}] Para que ocorra a sobrescrita, o método na subclasse deve possuir o mesmo nome e a mesma assinatura (parâmetros) do método na superclasse.
        \item[\textbf{( F )}] Um método definido como \texttt{final} na superclasse pode ser livremente sobrescrito... (\textit{Métodos final não podem ser sobrescritos}).
        \item[\textbf{( V )}] O uso da anotação \texttt{@Override} é opcional, mas ajuda o compilador...
        \item[\textbf{( V )}] A sobrescrita permite que uma subclasse forneça uma implementação específica...
        \item[\textbf{( F )}] Se um método for declarado como \texttt{static} na superclasse, ele pode ser sobrescrito... (\textit{Métodos estáticos são ocultados, não sobrescritos}).
    \end{enumerate}

    \item \textbf{Resposta:}\\
    \resposta{0.5cm}{
    a) O compilador baseia-se no tipo da referência (\texttt{Animal}) para verificar a validade das chamadas. Como \texttt{Animal} não possui o método \texttt{enterrarOsso()}, a chamada direta é negada por segurança de tipo. \\
    b) Ocorreria uma exceção do tipo \texttt{ClassCastException} em tempo de execução, pois um objeto do tipo \texttt{Gato} não pode ser convertido para \texttt{Cachorro}, já que não há relação de herança entre eles (são irmãos na hierarquia).
    }

    \item \textbf{Resposta: A}
    \begin{itemize}
        \item I: Correta. \texttt{final} em classes impede herança.
        \item II: Incorreta. \texttt{final} em métodos impede sobrescrita.
        \item III: Correta. \texttt{final} em atributos cria constantes (após inicialização).
        \item IV: Correta. Sem herança, não há sobrescrita.
    \end{itemize}

\end{enumerate}

\newpage
\begin{center}
    \rule{\textwidth}{0.4pt}
    \textbf{QUESTÕES PRÁTICAS}
\end{center}

\begin{enumerate}
    \item \textbf{Implementação de \texttt{ItemVestuario}:}\\
    \resposta{4cm}{
    \texttt{public abstract class ItemVestuario \{ \\
    \quad protected String marca; \\
    \quad protected double preco; \\
    \quad public static final String MOEDA = "R\$"; \\
    \quad public static int totalItens = 0; \\
    \quad \\
    \quad public abstract void exibirDetalhes(); \\
    \}}
    }

    \item \textbf{Implementação de \texttt{Camiseta}:}\\
    \resposta{6cm}{
    \texttt{public class Camiseta extends ItemVestuario \{ \\
    \quad private String tamanho; \\
    \quad public static final String TIPO = "Camiseta"; \\
    \quad \\
    \quad public Camiseta(String marca, double preco, String tamanho) \{ \\
    \quad\quad this.marca = marca; \\
    \quad\quad this.preco = preco; \\
    \quad\quad this.tamanho = tamanho; \\
    \quad\quad totalItens++; \\
    \quad \} \\
    \quad \\
    \quad @Override \\
    \quad public void exibirDetalhes() \{ \\
    \quad\quad System.out.println("Tipo: " + TIPO + " | Marca: " + marca + \\
    \quad\quad\quad " | Preço: " + MOEDA + preco + " | Tam: " + tamanho); \\
    \quad \} \\
    \}}
    }

    \item \textbf{Implementação de \texttt{Blusa}:}\\
    \resposta{6cm}{
    \texttt{public class Blusa extends ItemVestuario \{ \\
    \quad private String tecido; \\
    \quad public static final String TIPO = "Blusa"; \\
    \quad \\
    \quad public Blusa(String marca, double preco, String tecido) \{ \\
    \quad\quad this.marca = marca; \\
    \quad\quad this.preco = preco; \\
    \quad\quad this.tecido = tecido; \\
    \quad\quad totalItens++; \\
    \quad \} \\
    \quad \\
    \quad @Override \\
    \quad public void exibirDetalhes() \{ \\
    \quad\quad System.out.println("Tipo: " + TIPO + " | Marca: " + marca + \\
    \quad\quad\quad " | Preço: " + MOEDA + preco + " | Tecido: " + tecido); \\
    \quad \} \\
    \}}
    }

    \item \textbf{Downcasting e Explicação:}\\
    \resposta{3cm}{
    \texttt{((Camiseta) item).dobrarCamiseta();} \\
    Os parênteses externos são necessários devido à precedência de operadores. O operador de ponto (\texttt{.}) tem precedência maior que o operador de \textit{cast}. Sem os parênteses, o compilador tentaria primeiro executar o método no objeto do tipo \texttt{ItemVestuario} antes de realizar a conversão.
    }

    \item \textbf{SOLID e Classes Abstratas:}\\
    \resposta{3cm}{
    As classes abstratas são fundamentais para o Princípio do Aberto/Fechado (OCP), permitindo que o comportamento de um sistema seja estendido sem modificar o código existente, e para o Princípio da Inversão de Dependência (DIP), pois garantem que módulos de alto nível dependam de abstrações e não de implementações concretas, promovendo baixo acoplamento.
    }
\end{enumerate}

\end{document} 
